\mathplain
\section{Compared with the Direct Proof}\label{sec:comparison}

The proof of Rapoport et al.\ \cite{rapoport2017simple} is the most direct of the published proofs, and it helps to put the two side by side. This section adapts its invariant to store notation. Only the translation route is mechanized here, so the comparison is informal.

\subsection{Two invariants}

The direct proof keeps a source typing along the run. Its invariant is that the state is typed in an inert context, a context in which every binder has the type of a function or a recursive object type whose type members have exact bounds, as the type of a value does. The translation keeps a target typing instead. Its invariant is that some typed target state has the same erasure as the source state. Both proofs must recover a lambda at an application and a field at a projection, and the difference is where they recover it from.

\subsection{The lemmas, side by side}

The direct proof builds a ladder of typing judgments. Precise typing $\src{\Gamma \vdash_! x : T}$ starts at the declared type $\src{\Gamma(x)}$ and only unfolds recursive types and projects intersections, so it exposes the declarations that the context records. Tight typing $\src{\Gamma \vdash_\# x : T}$ restricts the selection rules to members read by precise typing with exact bounds:
\[
\src{\inferrule{\Gamma \vdash_! x : \dtyp{A}{U}{U}}{\Gamma \vdash_\# U \sub x.A \qquad \Gamma \vdash_\# x.A \sub U}} .
\]
Invertible typing separates eliminations from introductions so that induction recovers the precise type. The main theorem says that in an inert context general typing implies tight typing, and its main case, selection replacement, says that a tight typing $\src{\Gamma \vdash_\# x : \dtyp{A}{S}{T}}$ comes from a precise $\src{\dtyp{A}{U}{U}}$ with $\src{S \sub U \sub T}$. Each of these has a counterpart in \FCdot.

\begin{center}
\small
\begin{tabular}{@{}p{0.42\textwidth}p{0.5\textwidth}@{}}
\toprule
direct proof & \leavevmode\FCdot and the translation \\
\midrule
inert context & typed store, inert by construction: every location holds a lambda or a literal at its precise type \\
precise typing & the precise view of a stored literal \\
tight selection rules & the $\ev{member}$ case of the normalizer, which reads the definition of the root through the view \\
invertible typing & typed head forms, invertible by construction \\
general typing implies tight typing & the canonical forms theorem \\
selection replacement & the $\ev{member}$ case of the theorem \\
narrowing & substitution of a cast atom for a variable \\
transitivity in tight typing & composition of typed forms \\
canonical forms of values & Corollary~\ref{cor:inversion} \\
\bottomrule
\end{tabular}
\end{center}

\paragraph{Allocation.}
In the direct proof, a value $v$ reaches a binder $x : T$ and must be stored at an inert type. The proof recovers an inert type $P$ of $v$ with $\src{P \sub T}$, narrows the continuation from $x : T$ to $x : P$, and stores $v$ at $P$. The term does not change. In \FCdot the value is a literal under casts. The machine stores the literal at its precise type and substitutes $\ev{cast}\,x\,\weak{\bar{e}}$ for $x$ in the continuation. The term changes, and the typing of the new term is atom substitution, which is proven once, before anything about normal forms.

\paragraph{Application.}
In the direct proof, $f\,y$ with $f : \src{\dall{x}{S}{T}}$ inverts the typing of $f$ to find a lambda $\src{\lambda(x : S_0).\,t_0}$ in the store with $\src{S \sub S_0}$ and a body of type $T_0$ with $\src{T_0 \sub T}$ under $x : S$. It narrows the body to $x : S$, subsumes its result, and substitutes $y$. The proof reconstructs why the body meets the advertised type. In \FCdot the corresponding facts are the domain and codomain evidence of the function form that the chain of casts normalizes to, and the machine uses them: it casts the argument by $d$ and the result by $c$. The body is never narrowed, because the cast on the argument does that work at run time.

\paragraph{Projection.}
In the direct proof, $x.a$ with $\src{x : \dfld{a}{T}}$ inverts the typing of $x$ to find the field body and a type $S$ of it with $\src{S \sub T}$. In \FCdot a projection carries presence evidence, which canonical forms turn into a field of the literal at the root. The type of the field is the block name $\sel{x}{a}$, the field body is typed at that name, and the cast that the translation put around the projection converts the name to the declared type. No inversion of field subtyping is used.

\paragraph{Transitivity.}
Both proofs need a composition theorem. In the direct proof it is about tight subtyping derivations. In \FCdot it is Lemma~\ref{lem:combine}, about typed forms. It is an induction on the sizes of the two forms, with no measure on derivations. That works because the entries of object forms are normalized templates: composing two object forms substitutes entries into entries without consulting any view.

\subsection{What each route costs}

The direct proof proves preservation for DOT, which the translation does not. The translation proves safety only, through the target. The translation costs a second calculus, its checker, and the translation itself, and the development is larger than a direct proof of one calculus would need to be. What it buys is that the lemmas about evidence are lemmas about \FCdot, not about the derivations of one source calculus. Extending the source with paths, captures, or a new sort of member extends \FCdot with constructors and the normalizer with cases. The theorem keeps its statement, and preservation and progress keep their shape, though they gain cases. Section~\ref{sec:conclusion} says how far that has been tested.
